%%%%%%%%%%%%%%%%%%%%%%%%%%%%%%%%%%%%%%%%%%%%%%%%%%%%%%%%%%%%%%%%%%%%%%%%
% Section introduction
% Inclus depuis main.tex via %%%%%%%%%%%%%%%%%%%%%%%%%%%%%%%%%%%%%%%%%%%%%%%%%%%%%%%%%%%%%%%%%%%%%%%%
% Section introduction
% Inclus depuis main.tex via %%%%%%%%%%%%%%%%%%%%%%%%%%%%%%%%%%%%%%%%%%%%%%%%%%%%%%%%%%%%%%%%%%%%%%%%
% Section introduction
% Inclus depuis main.tex via %%%%%%%%%%%%%%%%%%%%%%%%%%%%%%%%%%%%%%%%%%%%%%%%%%%%%%%%%%%%%%%%%%%%%%%%
% Section introduction
% Inclus depuis main.tex via \input{chapters/introduction}
%%%%%%%%%%%%%%%%%%%%%%%%%%%%%%%%%%%%%%%%%%%%%%%%%%%%%%%%%%%%%%%%%%%%%%%%

%%%%%%%%%%%%%%%%%%%%%%%%%%%%%%%%%%%%%%%%%%%%%%%%%%%%%%%%%%%%%%%%%%%%%%%%

\subsection{Contexte : le TAL clinique au défi de la vie privée}

La généralisation des dossiers médicaux électroniques a transformé en profondeur la
documentation des soins. Environ 80~\% des données cliniques produites quotidiennement
se présentent sous forme de texte libre : comptes rendus d'hospitalisation, notes d'évolution,
comptes rendus de radiologie et évaluations infirmières. Ces narratifs capturent la finesse du
raisonnement clinique, la trajectoire du patient et les déterminants sociaux de santé que les
champs structurés ne savent pas représenter. Le traitement automatique des langues (TAL)
clinique est devenu la technologie clé pour exploiter cette richesse, avec des applications
allant de la classification diagnostique automatique à l'aide à la décision, en passant par
l'identification de cohortes pour la recherche~\citep{wang2018clinical, kreimeyer2017natural}.

Les progrès récents du domaine ont été portés par la mise à l'échelle des modèles de langue :
GatorTron~\citep{yang2022large}, entraîné sur \num{90} milliards de mots cliniques, ou encore Med-PaLM~2~\citep{singhal2023medpalm2}, qui
atteint \SI{86.5}{\percent} d'exactitude sur MedQA, illustrent un principe devenu dominant :
plus de données, de meilleurs modèles. Pourtant, dans le domaine médical, cet appétit en
données entre en collision frontale avec l'impératif de protection de la vie privée des patients.

\subsection{Tension fondamentale entre données et vie privée}

Les données de santé occupent une position privilégiée dans les cadres réglementaires
mondiaux. Le Règlement Général sur la Protection des Données (RGPD) en Europe~\citep{gdpr2016} et le
Health Insurance Portability and Accountability Act (HIPAA) aux États-Unis~\citep{hipaa1996} les classent
parmi les catégories les plus sensibles, soumises à des interdictions générales de traitement
assorties d'exceptions strictement encadrées. Les approches classiques de partage,
principalement l'anonymisation et la pseudonymisation, souffrent de limites structurelles :
une étude séminale~\citep{sweeney2002} a montré que \SI{87}{\percent} de la population américaine pouvait être
réidentifiée uniquement par la combinaison du sexe, de la date de naissance et du code
postal à cinq chiffres, informations couramment conservées dans les jeux de données
prétendus anonymisés. Plus fondamentalement, l'anonymisation agressive dégrade
précisément la richesse linguistique et clinique qui fait la valeur des textes pour le TAL.

Une seconde limite, souvent sous-estimée, concerne les modèles eux-mêmes : les réseaux
de neurones peuvent mémoriser et divulguer des exemples d'entraînement. Les attaques par
inférence d'appartenance~\citep{shokri2017membership} ou par extraction de données d'entraînement~\citep{carlini2021extracting} démontrent que
partager un modèle entraîné sur des textes cliniques sensibles peut présenter un risque
comparable à celui du partage direct des données sources. Il en résulte un paysage
fragmenté où les rares corpus publics (\mimic{}, i2b2/n2c2) servent de \emph{de facto}
standards, et où les outils véritablement utiles restent cantonnés à l'intérieur des
institutions qui les développent.

\subsection{Question de recherche et hypothèse des données synthétiques}

Les données synthétiques offrent une voie prometteuse : des documents générés
artificiellement qui imitent les propriétés statistiques et linguistiques des notes cliniques
réelles sans correspondre à aucun patient identifiable. Les avancées récentes en modélisation
générative, et en particulier les grands modèles de langue, rendent ce programme crédible.
La présente thèse investigue cette piste à travers la question suivante :
\begin{center}
\textit{Comment générer des documents cliniques synthétiques à la fois utiles pour
l'entraînement de modèles de TAL et sûrs au regard des exigences réglementaires ?}
\end{center}

Cette interrogation centrale se décline en quatre questions de recherche complémentaires :
\begin{enumerate}
  \item \textbf{Q1 (génération)} Comment tirer parti des grands modèles de langue pour
    produire des documents cliniques synthétiques linguistiquement réalistes et
    médicalement cohérents ?
  \item \textbf{Q2 (utilité aval)} Dans quelle mesure les documents synthétiques peuvent-ils
    servir de données d'entraînement efficaces pour les tâches de TAL clinique en aval, et
    à quelles conditions apportent-ils un gain par rapport aux seules données réelles ?
  \item \textbf{Q3 (annotation faible)} Les grands modèles de langue peuvent-ils jouer le
    rôle d'annotateurs efficaces pour les textes cliniques, notamment dans les contextes
    multilingues et faiblement dotés ?
  \item \textbf{Q4 (vie privée)} Quelles sont les implications de la génération synthétique
    sur la vie privée, et comment concevoir des approches qui minimisent les risques de
    réidentification tout en préservant l'utilité des données ?
\end{enumerate}

\subsection{Contributions et plan du résumé}

Cette thèse développe une approche dite \textit{fac-similé} : les documents synthétiques ne
sont pas générés depuis rien, mais ancrés dans des narratifs réels par l'intermédiaire de
mots-clés médicaux extraits automatiquement (codes diagnostiques, procédures, médicaments,
symptômes). Seuls ces mots-clés et des scores numériques de qualité traversent une
frontière de confidentialité explicite, ce qui permet de préserver la structure et la cohérence
cliniques sans jamais exposer le texte patient.

Les contributions s'articulent autour de quatre axes :
\begin{enumerate}
  \item une architecture de génération fac-similé reposant sur l'extraction de concepts
    \gls{umls} et un raffinement itératif par optimisation de préférence directe guidé par un
    scoreur contrastif (Section~\ref{sec:facsimile} et Section~\ref{sec:dpo}) ;
  \item une validation translinguistique sur le français clinique (données \gls{aphp}, CIM-10)
    démontrant l'indépendance de la méthodologie vis-à-vis de la langue et du système de
    codage (Section~\ref{sec:french}) ;
  \item une étude de distillation de grands modèles de langue pour l'annotation faible
    d'entités cliniques en cinq langues européennes, dont le basque faiblement doté
    (Section~\ref{sec:weak}) ;
  \item une analyse rigoureuse des risques de réidentification résiduels par attaques par
    liaison et de discrimination, accompagnée de stratégies d'atténuation ciblées
    (Section~\ref{sec:privacy}).
\end{enumerate}

Le présent résumé suit cette structure par contribution. L'état de l'art
(Section~\ref{sec:related}) précède la présentation des contributions, et une discussion
finale (Section~\ref{sec:conclusion}) reprend les réponses explicites aux quatre questions
de recherche et expose les perspectives.

%%%%%%%%%%%%%%%%%%%%%%%%%%%%%%%%%%%%%%%%%%%%%%%%%%%%%%%%%%%%%%%%%%%%%%%%

%%%%%%%%%%%%%%%%%%%%%%%%%%%%%%%%%%%%%%%%%%%%%%%%%%%%%%%%%%%%%%%%%%%%%%%%

\subsection{Contexte : le TAL clinique au défi de la vie privée}

La généralisation des dossiers médicaux électroniques a transformé en profondeur la
documentation des soins. Environ 80~\% des données cliniques produites quotidiennement
se présentent sous forme de texte libre : comptes rendus d'hospitalisation, notes d'évolution,
comptes rendus de radiologie et évaluations infirmières. Ces narratifs capturent la finesse du
raisonnement clinique, la trajectoire du patient et les déterminants sociaux de santé que les
champs structurés ne savent pas représenter. Le traitement automatique des langues (TAL)
clinique est devenu la technologie clé pour exploiter cette richesse, avec des applications
allant de la classification diagnostique automatique à l'aide à la décision, en passant par
l'identification de cohortes pour la recherche~\citep{wang2018clinical, kreimeyer2017natural}.

Les progrès récents du domaine ont été portés par la mise à l'échelle des modèles de langue :
GatorTron~\citep{yang2022large}, entraîné sur \num{90} milliards de mots cliniques, ou encore Med-PaLM~2~\citep{singhal2023medpalm2}, qui
atteint \SI{86.5}{\percent} d'exactitude sur MedQA, illustrent un principe devenu dominant :
plus de données, de meilleurs modèles. Pourtant, dans le domaine médical, cet appétit en
données entre en collision frontale avec l'impératif de protection de la vie privée des patients.

\subsection{Tension fondamentale entre données et vie privée}

Les données de santé occupent une position privilégiée dans les cadres réglementaires
mondiaux. Le Règlement Général sur la Protection des Données (RGPD) en Europe~\citep{gdpr2016} et le
Health Insurance Portability and Accountability Act (HIPAA) aux États-Unis~\citep{hipaa1996} les classent
parmi les catégories les plus sensibles, soumises à des interdictions générales de traitement
assorties d'exceptions strictement encadrées. Les approches classiques de partage,
principalement l'anonymisation et la pseudonymisation, souffrent de limites structurelles :
une étude séminale~\citep{sweeney2002} a montré que \SI{87}{\percent} de la population américaine pouvait être
réidentifiée uniquement par la combinaison du sexe, de la date de naissance et du code
postal à cinq chiffres, informations couramment conservées dans les jeux de données
prétendus anonymisés. Plus fondamentalement, l'anonymisation agressive dégrade
précisément la richesse linguistique et clinique qui fait la valeur des textes pour le TAL.

Une seconde limite, souvent sous-estimée, concerne les modèles eux-mêmes : les réseaux
de neurones peuvent mémoriser et divulguer des exemples d'entraînement. Les attaques par
inférence d'appartenance~\citep{shokri2017membership} ou par extraction de données d'entraînement~\citep{carlini2021extracting} démontrent que
partager un modèle entraîné sur des textes cliniques sensibles peut présenter un risque
comparable à celui du partage direct des données sources. Il en résulte un paysage
fragmenté où les rares corpus publics (\mimic{}, i2b2/n2c2) servent de \emph{de facto}
standards, et où les outils véritablement utiles restent cantonnés à l'intérieur des
institutions qui les développent.

\subsection{Question de recherche et hypothèse des données synthétiques}

Les données synthétiques offrent une voie prometteuse : des documents générés
artificiellement qui imitent les propriétés statistiques et linguistiques des notes cliniques
réelles sans correspondre à aucun patient identifiable. Les avancées récentes en modélisation
générative, et en particulier les grands modèles de langue, rendent ce programme crédible.
La présente thèse investigue cette piste à travers la question suivante :
\begin{center}
\textit{Comment générer des documents cliniques synthétiques à la fois utiles pour
l'entraînement de modèles de TAL et sûrs au regard des exigences réglementaires ?}
\end{center}

Cette interrogation centrale se décline en quatre questions de recherche complémentaires :
\begin{enumerate}
  \item \textbf{Q1 (génération)} Comment tirer parti des grands modèles de langue pour
    produire des documents cliniques synthétiques linguistiquement réalistes et
    médicalement cohérents ?
  \item \textbf{Q2 (utilité aval)} Dans quelle mesure les documents synthétiques peuvent-ils
    servir de données d'entraînement efficaces pour les tâches de TAL clinique en aval, et
    à quelles conditions apportent-ils un gain par rapport aux seules données réelles ?
  \item \textbf{Q3 (annotation faible)} Les grands modèles de langue peuvent-ils jouer le
    rôle d'annotateurs efficaces pour les textes cliniques, notamment dans les contextes
    multilingues et faiblement dotés ?
  \item \textbf{Q4 (vie privée)} Quelles sont les implications de la génération synthétique
    sur la vie privée, et comment concevoir des approches qui minimisent les risques de
    réidentification tout en préservant l'utilité des données ?
\end{enumerate}

\subsection{Contributions et plan du résumé}

Cette thèse développe une approche dite \textit{fac-similé} : les documents synthétiques ne
sont pas générés depuis rien, mais ancrés dans des narratifs réels par l'intermédiaire de
mots-clés médicaux extraits automatiquement (codes diagnostiques, procédures, médicaments,
symptômes). Seuls ces mots-clés et des scores numériques de qualité traversent une
frontière de confidentialité explicite, ce qui permet de préserver la structure et la cohérence
cliniques sans jamais exposer le texte patient.

Les contributions s'articulent autour de quatre axes :
\begin{enumerate}
  \item une architecture de génération fac-similé reposant sur l'extraction de concepts
    \gls{umls} et un raffinement itératif par optimisation de préférence directe guidé par un
    scoreur contrastif (Section~\ref{sec:facsimile} et Section~\ref{sec:dpo}) ;
  \item une validation translinguistique sur le français clinique (données \gls{aphp}, CIM-10)
    démontrant l'indépendance de la méthodologie vis-à-vis de la langue et du système de
    codage (Section~\ref{sec:french}) ;
  \item une étude de distillation de grands modèles de langue pour l'annotation faible
    d'entités cliniques en cinq langues européennes, dont le basque faiblement doté
    (Section~\ref{sec:weak}) ;
  \item une analyse rigoureuse des risques de réidentification résiduels par attaques par
    liaison et de discrimination, accompagnée de stratégies d'atténuation ciblées
    (Section~\ref{sec:privacy}).
\end{enumerate}

Le présent résumé suit cette structure par contribution. L'état de l'art
(Section~\ref{sec:related}) précède la présentation des contributions, et une discussion
finale (Section~\ref{sec:conclusion}) reprend les réponses explicites aux quatre questions
de recherche et expose les perspectives.

%%%%%%%%%%%%%%%%%%%%%%%%%%%%%%%%%%%%%%%%%%%%%%%%%%%%%%%%%%%%%%%%%%%%%%%%

%%%%%%%%%%%%%%%%%%%%%%%%%%%%%%%%%%%%%%%%%%%%%%%%%%%%%%%%%%%%%%%%%%%%%%%%

\subsection{Contexte : le TAL clinique au défi de la vie privée}

La généralisation des dossiers médicaux électroniques a transformé en profondeur la
documentation des soins. Environ 80~\% des données cliniques produites quotidiennement
se présentent sous forme de texte libre : comptes rendus d'hospitalisation, notes d'évolution,
comptes rendus de radiologie et évaluations infirmières. Ces narratifs capturent la finesse du
raisonnement clinique, la trajectoire du patient et les déterminants sociaux de santé que les
champs structurés ne savent pas représenter. Le traitement automatique des langues (TAL)
clinique est devenu la technologie clé pour exploiter cette richesse, avec des applications
allant de la classification diagnostique automatique à l'aide à la décision, en passant par
l'identification de cohortes pour la recherche~\citep{wang2018clinical, kreimeyer2017natural}.

Les progrès récents du domaine ont été portés par la mise à l'échelle des modèles de langue :
GatorTron~\citep{yang2022large}, entraîné sur \num{90} milliards de mots cliniques, ou encore Med-PaLM~2~\citep{singhal2023medpalm2}, qui
atteint \SI{86.5}{\percent} d'exactitude sur MedQA, illustrent un principe devenu dominant :
plus de données, de meilleurs modèles. Pourtant, dans le domaine médical, cet appétit en
données entre en collision frontale avec l'impératif de protection de la vie privée des patients.

\subsection{Tension fondamentale entre données et vie privée}

Les données de santé occupent une position privilégiée dans les cadres réglementaires
mondiaux. Le Règlement Général sur la Protection des Données (RGPD) en Europe~\citep{gdpr2016} et le
Health Insurance Portability and Accountability Act (HIPAA) aux États-Unis~\citep{hipaa1996} les classent
parmi les catégories les plus sensibles, soumises à des interdictions générales de traitement
assorties d'exceptions strictement encadrées. Les approches classiques de partage,
principalement l'anonymisation et la pseudonymisation, souffrent de limites structurelles :
une étude séminale~\citep{sweeney2002} a montré que \SI{87}{\percent} de la population américaine pouvait être
réidentifiée uniquement par la combinaison du sexe, de la date de naissance et du code
postal à cinq chiffres, informations couramment conservées dans les jeux de données
prétendus anonymisés. Plus fondamentalement, l'anonymisation agressive dégrade
précisément la richesse linguistique et clinique qui fait la valeur des textes pour le TAL.

Une seconde limite, souvent sous-estimée, concerne les modèles eux-mêmes : les réseaux
de neurones peuvent mémoriser et divulguer des exemples d'entraînement. Les attaques par
inférence d'appartenance~\citep{shokri2017membership} ou par extraction de données d'entraînement~\citep{carlini2021extracting} démontrent que
partager un modèle entraîné sur des textes cliniques sensibles peut présenter un risque
comparable à celui du partage direct des données sources. Il en résulte un paysage
fragmenté où les rares corpus publics (\mimic{}, i2b2/n2c2) servent de \emph{de facto}
standards, et où les outils véritablement utiles restent cantonnés à l'intérieur des
institutions qui les développent.

\subsection{Question de recherche et hypothèse des données synthétiques}

Les données synthétiques offrent une voie prometteuse : des documents générés
artificiellement qui imitent les propriétés statistiques et linguistiques des notes cliniques
réelles sans correspondre à aucun patient identifiable. Les avancées récentes en modélisation
générative, et en particulier les grands modèles de langue, rendent ce programme crédible.
La présente thèse investigue cette piste à travers la question suivante :
\begin{center}
\textit{Comment générer des documents cliniques synthétiques à la fois utiles pour
l'entraînement de modèles de TAL et sûrs au regard des exigences réglementaires ?}
\end{center}

Cette interrogation centrale se décline en quatre questions de recherche complémentaires :
\begin{enumerate}
  \item \textbf{Q1 (génération)} Comment tirer parti des grands modèles de langue pour
    produire des documents cliniques synthétiques linguistiquement réalistes et
    médicalement cohérents ?
  \item \textbf{Q2 (utilité aval)} Dans quelle mesure les documents synthétiques peuvent-ils
    servir de données d'entraînement efficaces pour les tâches de TAL clinique en aval, et
    à quelles conditions apportent-ils un gain par rapport aux seules données réelles ?
  \item \textbf{Q3 (annotation faible)} Les grands modèles de langue peuvent-ils jouer le
    rôle d'annotateurs efficaces pour les textes cliniques, notamment dans les contextes
    multilingues et faiblement dotés ?
  \item \textbf{Q4 (vie privée)} Quelles sont les implications de la génération synthétique
    sur la vie privée, et comment concevoir des approches qui minimisent les risques de
    réidentification tout en préservant l'utilité des données ?
\end{enumerate}

\subsection{Contributions et plan du résumé}

Cette thèse développe une approche dite \textit{fac-similé} : les documents synthétiques ne
sont pas générés depuis rien, mais ancrés dans des narratifs réels par l'intermédiaire de
mots-clés médicaux extraits automatiquement (codes diagnostiques, procédures, médicaments,
symptômes). Seuls ces mots-clés et des scores numériques de qualité traversent une
frontière de confidentialité explicite, ce qui permet de préserver la structure et la cohérence
cliniques sans jamais exposer le texte patient.

Les contributions s'articulent autour de quatre axes :
\begin{enumerate}
  \item une architecture de génération fac-similé reposant sur l'extraction de concepts
    \gls{umls} et un raffinement itératif par optimisation de préférence directe guidé par un
    scoreur contrastif (Section~\ref{sec:facsimile} et Section~\ref{sec:dpo}) ;
  \item une validation translinguistique sur le français clinique (données \gls{aphp}, CIM-10)
    démontrant l'indépendance de la méthodologie vis-à-vis de la langue et du système de
    codage (Section~\ref{sec:french}) ;
  \item une étude de distillation de grands modèles de langue pour l'annotation faible
    d'entités cliniques en cinq langues européennes, dont le basque faiblement doté
    (Section~\ref{sec:weak}) ;
  \item une analyse rigoureuse des risques de réidentification résiduels par attaques par
    liaison et de discrimination, accompagnée de stratégies d'atténuation ciblées
    (Section~\ref{sec:privacy}).
\end{enumerate}

Le présent résumé suit cette structure par contribution. L'état de l'art
(Section~\ref{sec:related}) précède la présentation des contributions, et une discussion
finale (Section~\ref{sec:conclusion}) reprend les réponses explicites aux quatre questions
de recherche et expose les perspectives.

%%%%%%%%%%%%%%%%%%%%%%%%%%%%%%%%%%%%%%%%%%%%%%%%%%%%%%%%%%%%%%%%%%%%%%%%

%%%%%%%%%%%%%%%%%%%%%%%%%%%%%%%%%%%%%%%%%%%%%%%%%%%%%%%%%%%%%%%%%%%%%%%%

\subsection{Contexte : le TAL clinique au défi de la vie privée}

La généralisation des dossiers médicaux électroniques a transformé en profondeur la
documentation des soins. Environ 80~\% des données cliniques produites quotidiennement
se présentent sous forme de texte libre : comptes rendus d'hospitalisation, notes d'évolution,
comptes rendus de radiologie et évaluations infirmières. Ces narratifs capturent la finesse du
raisonnement clinique, la trajectoire du patient et les déterminants sociaux de santé que les
champs structurés ne savent pas représenter. Le traitement automatique des langues (TAL)
clinique est devenu la technologie clé pour exploiter cette richesse, avec des applications
allant de la classification diagnostique automatique à l'aide à la décision, en passant par
l'identification de cohortes pour la recherche~\citep{wang2018clinical, kreimeyer2017natural}.

Les progrès récents du domaine ont été portés par la mise à l'échelle des modèles de langue :
GatorTron~\citep{yang2022large}, entraîné sur \num{90} milliards de mots cliniques, ou encore Med-PaLM~2~\citep{singhal2023medpalm2}, qui
atteint \SI{86.5}{\percent} d'exactitude sur MedQA, illustrent un principe devenu dominant :
plus de données, de meilleurs modèles. Pourtant, dans le domaine médical, cet appétit en
données entre en collision frontale avec l'impératif de protection de la vie privée des patients.

\subsection{Tension fondamentale entre données et vie privée}

Les données de santé occupent une position privilégiée dans les cadres réglementaires
mondiaux. Le Règlement Général sur la Protection des Données (RGPD) en Europe~\citep{gdpr2016} et le
Health Insurance Portability and Accountability Act (HIPAA) aux États-Unis~\citep{hipaa1996} les classent
parmi les catégories les plus sensibles, soumises à des interdictions générales de traitement
assorties d'exceptions strictement encadrées. Les approches classiques de partage,
principalement l'anonymisation et la pseudonymisation, souffrent de limites structurelles :
une étude séminale~\citep{sweeney2002} a montré que \SI{87}{\percent} de la population américaine pouvait être
réidentifiée uniquement par la combinaison du sexe, de la date de naissance et du code
postal à cinq chiffres, informations couramment conservées dans les jeux de données
prétendus anonymisés. Plus fondamentalement, l'anonymisation agressive dégrade
précisément la richesse linguistique et clinique qui fait la valeur des textes pour le TAL.

Une seconde limite, souvent sous-estimée, concerne les modèles eux-mêmes : les réseaux
de neurones peuvent mémoriser et divulguer des exemples d'entraînement. Les attaques par
inférence d'appartenance~\citep{shokri2017membership} ou par extraction de données d'entraînement~\citep{carlini2021extracting} démontrent que
partager un modèle entraîné sur des textes cliniques sensibles peut présenter un risque
comparable à celui du partage direct des données sources. Il en résulte un paysage
fragmenté où les rares corpus publics (\mimic{}, i2b2/n2c2) servent de \emph{de facto}
standards, et où les outils véritablement utiles restent cantonnés à l'intérieur des
institutions qui les développent.

\subsection{Question de recherche et hypothèse des données synthétiques}

Les données synthétiques offrent une voie prometteuse : des documents générés
artificiellement qui imitent les propriétés statistiques et linguistiques des notes cliniques
réelles sans correspondre à aucun patient identifiable. Les avancées récentes en modélisation
générative, et en particulier les grands modèles de langue, rendent ce programme crédible.
La présente thèse investigue cette piste à travers la question suivante :
\begin{center}
\textit{Comment générer des documents cliniques synthétiques à la fois utiles pour
l'entraînement de modèles de TAL et sûrs au regard des exigences réglementaires ?}
\end{center}

Cette interrogation centrale se décline en quatre questions de recherche complémentaires :
\begin{enumerate}
  \item \textbf{Q1 (génération)} Comment tirer parti des grands modèles de langue pour
    produire des documents cliniques synthétiques linguistiquement réalistes et
    médicalement cohérents ?
  \item \textbf{Q2 (utilité aval)} Dans quelle mesure les documents synthétiques peuvent-ils
    servir de données d'entraînement efficaces pour les tâches de TAL clinique en aval, et
    à quelles conditions apportent-ils un gain par rapport aux seules données réelles ?
  \item \textbf{Q3 (annotation faible)} Les grands modèles de langue peuvent-ils jouer le
    rôle d'annotateurs efficaces pour les textes cliniques, notamment dans les contextes
    multilingues et faiblement dotés ?
  \item \textbf{Q4 (vie privée)} Quelles sont les implications de la génération synthétique
    sur la vie privée, et comment concevoir des approches qui minimisent les risques de
    réidentification tout en préservant l'utilité des données ?
\end{enumerate}

\subsection{Contributions et plan du résumé}

Cette thèse développe une approche dite \textit{fac-similé} : les documents synthétiques ne
sont pas générés depuis rien, mais ancrés dans des narratifs réels par l'intermédiaire de
mots-clés médicaux extraits automatiquement (codes diagnostiques, procédures, médicaments,
symptômes). Seuls ces mots-clés et des scores numériques de qualité traversent une
frontière de confidentialité explicite, ce qui permet de préserver la structure et la cohérence
cliniques sans jamais exposer le texte patient.

Les contributions s'articulent autour de quatre axes :
\begin{enumerate}
  \item une architecture de génération fac-similé reposant sur l'extraction de concepts
    \gls{umls} et un raffinement itératif par optimisation de préférence directe guidé par un
    scoreur contrastif (Section~\ref{sec:facsimile} et Section~\ref{sec:dpo}) ;
  \item une validation translinguistique sur le français clinique (données \gls{aphp}, CIM-10)
    démontrant l'indépendance de la méthodologie vis-à-vis de la langue et du système de
    codage (Section~\ref{sec:french}) ;
  \item une étude de distillation de grands modèles de langue pour l'annotation faible
    d'entités cliniques en cinq langues européennes, dont le basque faiblement doté
    (Section~\ref{sec:weak}) ;
  \item une analyse rigoureuse des risques de réidentification résiduels par attaques par
    liaison et de discrimination, accompagnée de stratégies d'atténuation ciblées
    (Section~\ref{sec:privacy}).
\end{enumerate}

Le présent résumé suit cette structure par contribution. L'état de l'art
(Section~\ref{sec:related}) précède la présentation des contributions, et une discussion
finale (Section~\ref{sec:conclusion}) reprend les réponses explicites aux quatre questions
de recherche et expose les perspectives.
